% titel: Blatt 7, A7: Molmasse und Dichte eines fluorierten Methans (CFxH4-x)
% Teil II: der durchgerechnete Weg.

\subsection*{Gegeben und gesucht}
$\rho = 8.0\ \mathrm{g\,L^{-1}}$ bei $p=2.81$ atm und $T=300$ K.
Gesucht: a) Molmasse, b) Lewis-Formel (nur C,H,F), c) Dichte bei $1.00$ atm
und $298$ K.

\subsection*{Ansatz}
$M=\rho RT/p$. Durch Abgleich mit $M(\mathrm{CF}_x\mathrm{H}_{4-x})$
$x$ finden. Bei c) $\rho\propto p/T$ nutzen.

\subsection*{Rechnung}
\textbf{a)} $M=\rho RT/p$:
\[
  M = \frac{8.0\cdot 0.0821\cdot 300}{2.81}
    = \frac{197.0}{2.81} = 70.1\ \mathrm{g\,mol^{-1}} .
\]
Abgleich $M(x)=12.01+19.00x+1.008(4-x)$:
\begin{itemize}
  \item $x=1$: $M=12.01+19.00+3.02 = 34.0$;
  \item $x=2$: $M=12.01+38.00+2.02 = 52.0$;
  \item $x=3$: $M=12.01+57.00+1.01 = 70.0$ $\checkmark$.
\end{itemize}
$\Rightarrow x=3$, Verbindung ist \textbf{Trifluormethan} $\mathrm{CHF_3}$
($= \mathrm{CF_3H}$).

\textbf{b)} Lewis: tetraedrisches C mit 3 C--F-Einfachbindungen und
1 C--H-Einfachbindung; C tr\"agt keine freien Paare, keine Doppelbindung.
\[
  \mathrm{H{-}C(F)_3}\quad\text{(tetraedrisch)}
\]

\textbf{c)} Bei $p_2=1.00$ atm, $T_2=298$ K (isotrop skaliert:
$M$ gleich):
\[
  \rho_2 = \rho_1\,\frac{p_2}{p_1}\cdot\frac{T_1}{T_2}
  = 8.0\cdot\frac{1.00}{2.81}\cdot\frac{300}{298}
  = 2.87\ \mathrm{g\,L^{-1}} .
\]

\begin{ergebnis}
a) $M=70.1\ \mathrm{g\,mol^{-1}}$ $\Rightarrow$ $\mathrm{CHF_3}$
(Trifluormethan). b) tetraedrisches C mit 3 C--F und 1 C--H.
c) $\rho=2.87\ \mathrm{g\,L^{-1}}$.
\end{ergebnis}

\begin{kontrolle}
C + 3 F + H: $12.01+3\times 19.00+1.01 = 70.0\ \mathrm{g\,mol^{-1}}$ deckt
sich mit der gemessenen $70.1$. Druckabfall von $2.81\to1$ atm und leicht
tiefere Temperatur senken die Dichte wie erwartet auf $2.87$ g/L
(und $M$ bleibt konstant).
\end{kontrolle}

\begin{merksatz}
Dichte des idealen Gases: $\rho=pM/(RT)$. Aus der Molmasse die
Summenformel bestimmen (ganzzahliges $x$); die Dichte skaliert mit
$p/T$ bei konstanter Molmasse.
\end{merksatz}
